\documentclass[11pt,a4paper]{article}
\usepackage[utf8]{inputenc}
\usepackage[T1]{fontenc}
\usepackage{lmodern,amsmath,amssymb,amsthm,mathtools,array,booktabs,enumitem,microtype}
\usepackage[margin=1in]{geometry}
\usepackage{natbib}
\usepackage{hyperref}
\usepackage{xr-hyper}
\externaldocument[PH-]{01a_graphic_equalisation_physics_grounding}

\newtheorem{definition}{Definition}[section]
\newtheorem{proposition}[definition]{Proposition}
\newtheorem{hypothesis}[definition]{Hypothesis}
\newtheorem{condition}[definition]{Condition}
\theoremstyle{remark}
\newtheorem{remark}[definition]{Remark}

\newcommand{\Cont}{\mathfrak C}
\newcommand{\Alt}{\mathcal A}
\newcommand{\State}{\mathcal X}
\newcommand{\Hist}{\mathcal H}
\newcommand{\Obs}{\mathcal O}

\title{The Law of Graphic Equalisation --- Passive Continuation Grounding:\\
\large Momentum, Gravity, and the Cost of Alteration}
\author{Matthew Parslow\\Independent Researcher}
\date{\textbf{Draft \texorpdfstring{\textemdash}{---}\ \PaperDraft}\\[2pt]\PaperDisplayDate \texorpdfstring{\textbullet}{\textperiodcentered}\ document revision \PaperRevision}

\usepackage{fancyhdr}
\pagestyle{fancy}
\fancyhf{}
\fancyfoot[C]{\thepage}
\fancyfoot[R]{\scriptsize\PaperFooterStamp}
\renewcommand{\headrulewidth}{0pt}
\fancypagestyle{plain}{\fancyhf{}\fancyfoot[C]{\thepage}%
  \fancyfoot[R]{\scriptsize\PaperFooterStamp}\renewcommand{\headrulewidth}{0pt}}
% GENERATED by tools/build_paper_metadata.py from paper_metadata.json.
% DO NOT EDIT: edit paper_metadata.json, then regenerate this file.
\providecommand{\PaperDraft}{draft-2026-09-06.1}
\providecommand{\PaperDisplayDate}{6 September 2026}
\expandafter\def\csname PaperRevision@00_is\endcsname{129}
\expandafter\def\csname PaperRevision@01_law_of_graphic_equalisation_theory_referenced\endcsname{294}
\expandafter\def\csname PaperRevision@01a_graphic_equalisation_physics_grounding\endcsname{47}
\expandafter\def\csname PaperRevision@01b_graphic_equalisation_mind_and_self_grounding\endcsname{47}
\expandafter\def\csname PaperRevision@01c_graphic_equalisation_cosmology_grounding\endcsname{11}
\expandafter\def\csname PaperRevision@01d_graphic_equalisation_born_recovery\endcsname{6}
\expandafter\def\csname PaperRevision@01e_graphic_equalisation_communication_grounding\endcsname{6}
\expandafter\def\csname PaperRevision@01f_graphic_equalisation_qualia_and_consciousness_grounding\endcsname{3}
\expandafter\def\csname PaperRevision@01g_graphic_equalisation_total_perspective_vortex_grounding\endcsname{4}
\expandafter\def\csname PaperRevision@01h_graphic_equalisation_information_gravity_quantum_reconciliation\endcsname{5}
\expandafter\def\csname PaperRevision@01i_graphic_equalisation_passive_continuation_momentum_gravity\endcsname{12}
\expandafter\def\csname PaperRevision@01s_graphic_equalisation_statement\endcsname{294}
\expandafter\def\csname PaperRevision@02_seldon_manifolds_theory_referenced\endcsname{31}
\expandafter\def\csname PaperRevision@02a_seldon_manifolds_psychological_grounding\endcsname{15}
\expandafter\def\csname PaperRevision@02b_seldon_manifolds_social_and_psychological_application\endcsname{7}
\expandafter\def\csname PaperRevision@02c_seldon_manifolds_decision_making\endcsname{11}
\expandafter\def\csname PaperRevision@03_law_of_narrative_momentum_theory_referenced\endcsname{32}
\expandafter\def\csname PaperRevision@04_narrativium_theory_referenced\endcsname{23}
\expandafter\def\csname PaperRevision@05_frames_relation_and_the_empty_limit\endcsname{11}
\expandafter\def\csname PaperRevision@06_assessing_an_account_of_this_kind\endcsname{29}
\expandafter\def\csname PaperRevision@07_causal_attitude_networks_on_the_substrate\endcsname{4}
\expandafter\def\csname PaperRevision@08_certified_answers_as_an_audit\endcsname{4}
\expandafter\def\csname PaperRevision@09_adaptability\endcsname{94}
\expandafter\def\csname PaperRevision@09a_adaptability_implementation\endcsname{7}
\expandafter\def\csname PaperRevision@10_applied_graphic_equalisation\endcsname{7}
\expandafter\def\csname PaperRevision@10a_red_queen_race_civilisational_rate_of_change_llms_biological_timescales_anthropocene_transition\endcsname{13}
\expandafter\def\csname PaperRevision@11_relational_information_ethics\endcsname{16}
\expandafter\def\csname PaperRevision@A1_constructive_graphic_equalisation_implementation_referenced\endcsname{8}
\expandafter\def\csname PaperGenerated@01s_graphic_equalisation_statement\endcsname{1}
\edef\PaperMetadataStem{\jobname}
\expandafter\ifx\csname PaperRevision@\PaperMetadataStem\endcsname\relax
  \PackageError{paper_metadata}{No metadata for \PaperMetadataStem}{Add the paper stem to paper_metadata.json.}
  \def\PaperRevision{?}
\else
  \expandafter\let\expandafter\PaperRevision\csname PaperRevision@\PaperMetadataStem\endcsname
\fi
\newcommand{\PaperDateLine}{\textbf{Draft \texorpdfstring{\textemdash}{---}\ \PaperDraft}\\[2pt]\PaperDisplayDate\ \texorpdfstring{\textbullet}{\textperiodcentered}\ document revision \PaperRevision}
\newcommand{\PaperStatementDateLine}{\textbf{Draft \texorpdfstring{\textemdash}{---}\ \PaperDraft}\\[2pt]\PaperDisplayDate\ \texorpdfstring{\textbullet}{\textperiodcentered}\ statement}
\expandafter\ifx\csname PaperGenerated@\PaperMetadataStem\endcsname\relax
  \newcommand{\PaperFooterStamp}{\textsc{\PaperDraft\,r\PaperRevision}}
\else
  \newcommand{\PaperFooterStamp}{\textsc{\PaperDraft\ / statement / from long r\PaperRevision}}
\fi
\newcommand{\PaperNoticeStamp}{\textbf{DRAFT --- \PaperDraft\ r\PaperRevision} (\PaperDisplayDate).}


\begin{document}
\maketitle
\begin{center}
\fbox{\parbox{0.92\textwidth}{\small
\PaperNoticeStamp\ This is Paper 01i, a formal companion to the Graphic
Equalisation, physics-grounding, total-perspective, information-gravity, and
Narrative Momentum papers. It formalises a research and derivation ledger, not
a new mechanics or gravitational theory. Its sole positive internal result is
class-relative: reuse of declared carried structure has zero \emph{additional}
alteration cost. The paper derives no physical momentum, four-momentum,
metric, connection, Einstein equation, stress--energy source, equivalence
principle, momentum-conservation law, geodesic equation, or quantum-gravity
result. Flat transport and load-shaped free propagation are recovery targets
for a concrete realization with stated equations, limits, and measurements.}}
\end{center}
\vspace{1em}

\begin{abstract}
Passive continuation is defined here as admissible reuse of presently operative
structure with no additional alteration cost in a declared grounding class. It
is not no activity, no time, no interaction at every grain, or no support cost
in another frame. This strict distinction permits a comparison with two
established targets without identifying their physical roles: momentum is not a
force, and general-relativistic free fall is not a locally felt gravitational
force on an ideal test body. In a flat regime a realization must establish
measured four-momentum transport, $dp^\mu/d\tau=0$, from a state map,
dynamics, symmetry, and isolation conditions. In a load-shaped regime it must
establish $u^\nu\nabla_\nu u^\mu=0$ from a Lorentzian metric and its
connection, a proper-time observable, and matter coupling, while separately supplying a
source law. Neither target follows from the inherited project vocabulary. The
conditional reconciliation investigated here is therefore exact: flat inertial
transport and load-shaped geodesic motion are two physical readings of one
passive-continuation law only if one concrete realization recovers their
equations, source/exchange accounting, universal-coupling limits, and stated
observables. A four-way ledger records what is inherited, hypothesised, still
owed as a physical bridge, and constrained by experiment.
\end{abstract}

\section{Scope and terminology}\label{sec:scope}
The question is not whether momentum and gravity are one force. They are not.
Physical momentum is a state/charge-like quantity; force is a specified rate of
its transfer or change. In general relativity (GR), an ideal freely falling test
body has zero proper acceleration, while a supported body records a
nongravitational support force. Curvature and tidal effects remain physically
measurable \citep{einstein1916,gerochjang1975,will2014}.

The question is narrower. Can one implemented law have a flat reading in which
carried physical momentum selects inertial transport and a load-shaped reading
in which geometry selects free propagation? This paper fixes the conditions
under which that question becomes a physics claim. It supplies neither the
implementation nor an inference that the conditions are already met.

The inherited corpus uses frame-indexed graph, admissibility, operative
history, and alteration-cost vocabulary. Its physics companion offers a
change-budget comparison; its Narrative Momentum companion distinguishes
preservation from alteration; its information-gravity companion requires typed
sources, geometry, clocks, and observables before a physical reconciliation
claim. The Stage-1 research record and Stage-2 derivation ledger for this paper
show that none provides the missing physical maps. Momentum and gravity are
historically distinguishable readings of a presently inseparable passive
continuation: a present in-frame record can measure a declared local relation,
but it cannot by itself attribute why the continuation is carried. That
attribution requires a history-bearing inference record with stated support,
source, and comparison conditions. It does not make momentum, gravity, inertia,
curvature, or force identical.

\begin{definition}[Declared continuation situation]\label{def:situation}
A frame-indexed continuation situation is
\[
 \Cont_t^W=(W,G_t,\Gamma_t,\mathcal P_t,S_t,\Alt,\kappa,\mathcal R_t,\Obs_t),
\]
where $W$ is the frame, $G_t$ the represented graph, $\Gamma_t$ its
admissibility grammar, $\mathcal P_t$ its admissible continuations, $S_t$ the
presently operative carried structure when represented, $\Alt$ a declared
alteration class, $\kappa$ its ordered cost domain, $\mathcal R_t$ relevant
support, and $\Obs_t$ its record relation. This is project-level typed
vocabulary; it is not a phase space, action, metric, clock, or source tensor.
\end{definition}
An alteration class must declare its support, equivalence relation on
alterations, comparison currency $\kappa$, realised reference condition,
resource budget, admissibility and measurement procedure. If any required
component is unavailable, the corresponding class cost is undefined, not
silently zero. The zero in the reuse convention belongs only to a fully
instantiated class.



\begin{definition}[History-bearing attribution record]\label{def:history}
For an interval $[t_0,t]$, a \emph{history-bearing attribution record} is
\[
 \Hist^W_{[t_0,t]}=
 \langle (\Obs_s,S_s,\mathcal R_s)_{t_0\leq s\leq t},\ I^W,\ \Lambda^W\rangle,
\]
where the retained sequence may be compressed but must preserve the stated
changes, supports, and comparison relations relevant to the typed attribution
map
\[
 I^W:\bigl((\Obs_s,S_s,\mathcal R_s)_{t_0\leq s\leq t},\Lambda^W\bigr)
 \longrightarrow \mathfrak A^W_{\rm phys}.
\]
The codomain $\mathfrak A^W_{\rm phys}$ is a declared set of typed physical
attributions (for example, momentum transport, load-shaped continuation,
curvature, or source claims), with their required observable and comparison
models. $\Lambda^W$ states the record's grain, losses, uncertainty, and domain
of use. The record is neither the total past nor a view from nowhere. It is the
minimum kind of record from which W can attempt to attribute a present passive
continuation to carried momentum-like history, load-shaped continuation
conditions, or both.
\end{definition}

\section{Inherited free-continuation and change-load claims}\label{sec:inherited}
Where $\kappa$ is additive, write
\[
 C_{\Alt}(a\mid S_t)=C_{0,\Alt}(a)+\Delta C_{\mathrm{alt},\Alt}(S_t,a).
\]
The inherited distinction licenses the following definition and no stronger
physical conclusion.

\begin{definition}[Passive continuation]\label{def:pc}
An admissible continuation $a\in\mathcal P_t$ is \emph{passive within $\Alt$}
when it reuses the declared carried structure and
\[
 \Delta C_{\mathrm{alt},\Alt}(S_t,a)=0.
\]
For an ordinal cost domain the valid statement is only the corresponding
no-extra-alteration ranking. Passive continuation means no \emph{additional}
alteration cost in $\Alt$; it does not mean zero total activity, zero total
cost, zero elapsed or inferred time, zero environmental support, or zero
interaction at every descriptive grain.
\end{definition}

\begin{proposition}[Internal passive-reuse result]\label{prop:reuse}
Given Definition~\ref{def:pc}, an admissible reuse continuation has zero
additional alteration cost in its declared class.
\end{proposition}
\begin{proof}
The conclusion is the defining cost assignment for reuse in that class. It is
not a conservation theorem because no physical state, unit, dynamics, symmetry,
or boundary current occurs in the premises.
\end{proof}

\begin{proposition}[History required for physical attribution]\label{prop:history}
A present passive-continuation relation does not determine a unique attribution
to physical momentum or load-shaped gravity. Any such attribution requires a
history-bearing record $\Hist^W_{[t_0,t]}$ and a declared inference map from
that record to the relevant physical observables.
\end{proposition}
\begin{proof}
Passive continuation classifies whether the present carried structure is reused;
it contains neither a physical momentum map nor a source/geometry map. A present
record can therefore be compatible with more than one physical realization. The
record in Definition~\ref{def:history} supplies the retained changes, supports,
and comparisons needed to test a proposed factorization, but still does not
supply the required physical map. Thus history is necessary for attribution
within this programme, though not sufficient without the declared inference
map.
\end{proof}

\begin{remark}[What in-frame inseparability does and does not say]
An unspecified in-frame record supplies no physical interpretation. With a
declared local accelerometer observable, a frame can determine proper
acceleration and distinguish free from supported motion in its stated domain.
That result does not differentiate momentum and gravity as explanatory factors:
both can remain readings of the present passive continuation. Their attribution
requires the history-bearing record of Definition~\ref{def:history}, together
with stated source, support, and comparison models. A local accelerometer cannot
thereby infer curvature or source structure, or distinguish a homogeneous
gravitational description from uniform acceleration without the specified
extended, source, or cross-frame inference. The limitation is therefore local
and operational: extended tidal measurements, support-force readings, source
modelling, and cross-frame comparison contribute to the required historical
inference. In particular, local inertial coordinates can remove connection
coefficients at one event, not curvature across a region \citep{einstein1916}.
\end{remark}

\section{Standard physical targets}\label{sec:targets}
Newtonian mechanics distinguishes momentum from force:
\[
 \mathbf p=m\mathbf v,\qquad \mathbf F_{\rm net}=\frac{d\mathbf p}{dt}.
\]
In an inertial frame, a free particle has constant momentum. For a closed
system, total momentum conservation additionally requires an appropriate action
or balance law, translation symmetry, and boundary-flux accounting; it does not
follow from persistence language alone \citep{newton1687,noether1918}.

Given a Lorentzian metric, a proper-time functional, and the point-particle
action $S_{\rm pp}=-mc\int d\tau$, the GR free-test-body target is
\[
 u^\nu\nabla_\nu u^\mu=0,
 \qquad u^\mu=\frac{dx^\mu}{d\tau}.
\]
It is equivalent in coordinates to the familiar connection equation. It is not
a field equation and does not identify the source of the metric. A supported
body has nonzero proper acceleration; a free body can have coordinate
acceleration without locally feeling a gravitational force
\citep{einstein1916,gerochjang1975}. Curvature appears operationally through,
for example, geodesic deviation,
\[
 \frac{D^2\xi^\mu}{D\tau^2}=-R^\mu{}_{\nu\rho\sigma}u^\nu\xi^\rho u^\sigma.
\]

Universality of free fall, local Lorentz and position behaviour, redshift,
light propagation, and post-Newtonian observables constrain any proposed
coupling \citep{will2014,wagner2012,touboul2022,bertotti2003}. They do not
show that inertial and gravitational concepts or mechanisms are identical.

\section{Formal passive-continuation definition}\label{sec:formal}
Definition~\ref{def:pc} is class-relative by construction. A containing frame
can price apparatus support or an environmental interaction even where a
mechanical-direction class assigns no additional alteration cost to reuse. No
sum between classes is licensed without a typed conversion. The definition
therefore has no physical units unless a realization supplies them.

\begin{proposition}[Persistence is not propulsion]\label{prop:notpropulsion}
Definition~\ref{def:pc} supplies no force, physical momentum, time parameter,
or equation of motion.
\end{proposition}
\begin{proof}
The definition fixes a relation in an ordered cost class. It specifies neither
a physical state space nor a map to momentum, nor a derivative, symmetry, or
interaction law. Each is required to state the listed physical objects.
\end{proof}

The strongest inherited result is thus structural: present operative structure
can be reused without an added alteration of that structure in its declared
class. Calling this a momentum-like reading is permissible only as a structural
analogy. Calling it physical momentum, force, energy, action, mass, or
stress--energy requires a separate map with units and measurements.

\section{Narrative residual-gradient derivation contract}\label{sec:informationforce}
This section makes the proposed direction of grounding explicit: physical gravity
is to be recovered from a realization of present operative Narrative, not used
to define the general continuation principle. The following is a new conditional
contract. It does not assert that the current corpus already supplies physical
gravity, a force in Newtons, or a source equation.

\begin{hypothesis}[H-IF: residual-gradient realization]\label{hyp:informationforce}
A candidate realization supplies a differential information-state space
$\mathcal Q_I$, a preparation/readout map
\[
 Z:\bigl(S_t,\Hist^W_{[t_0,t]},\mathcal R_t,\Obs_t\bigr)
 \longrightarrow q\in\mathcal Q_I,
\]
a physical clock/accelerometer bridge
\[
 \Theta:(q,\Obs_t)\longrightarrow(\tau,\mathcal A_{\rm acc}),
\]
where $\tau$ is an operational physical-clock reading and
$\mathcal A_{\rm acc}$ is a declared proper-acceleration/support observable,
a residual functional $\mathcal R_I:\mathcal Q_I\to K_I$, and a declared
ordered or differentiable value structure on $K_I$. Where its derivative is
defined, the \emph{residual gradient} is the residual derivative
\[
 \boxed{\mathcal F_I(q):=-D\mathcal R_I(q)\in T_q^*\mathcal Q_I.}\tag{IF1}
\]
A response map $B_q:T_q^*\mathcal Q_I\to T_q\mathcal Q_I$, if used to select
an update direction, must be supplied separately. Thus the residual gradient is
intrinsic to the declared information realization: it is the covector of
unresolved balance, not a metaphorical name for a later mechanical force.
\end{hypothesis}

H-IF has three mandatory physical translation stages. First, it must supply
measurement-preserving physical state and clock maps. Second, its Newton and
covariant force readouts must be equations, not identifications:
\[
\begin{aligned}
 K_{\rm N}&:(q,\mathcal F_I,\Theta)\longrightarrow\mathbb R^3,
 & [K_{\rm N}]&=\mathrm{kg\,m\,s^{-2}},
 &K_{\rm N}&=\frac{d\mathbf p}{dt},\\
 K_4&:(q,\mathcal F_I,\Theta)\longrightarrow T^*M,
 &K_4&=\frac{Dp_\mu}{d\tau}
 &&\text{in its declared relativistic convention}.
\end{aligned}\tag{IF2}
\]
Here $\Theta$ supplies the typed clock and accelerometer/readout conditions
above. The first line is the required map by which a newton becomes a
translation/readout unit for the candidate Narrative force; it does not give a
Narrative magnitude SI units by definition. Third, the realization must
independently close the existing source and geometry maps
\[
\begin{aligned}
 L&:(q,\mathcal F_I,\Theta)\longrightarrow T^{\rm eff}_{\mu\nu},\\
 \Phi&:(q,\mathcal F_I,\Theta)\longrightarrow\\[-2pt]
 &\qquad(M,g_{\mu\nu},\nabla,\Obs_{\rm phys}).
\end{aligned}\tag{IF3}
\]
Each map must state its tensor type and physical dimensions. It must also
specify an exchange law, matter coupling, and controlled error. A
force/exchange readout is not itself a rank-two gravitational source, and
neither source nor geometry may be assumed in order to define
$\mathcal R_I$.

\begin{proposition}[Conditional gravity-derivation target]\label{prop:ifrecovery}
If one realization satisfies H-IF, closes the maps in (IF2)--(IF3), and proves
F1--F6, IG1--IG8, H-PC, history-bearing attribution, universal coupling, and the
stated weak-field and precision limits, then that realization derives its
claimed physical gravitational regime from the residual-gradient construction.
In particular, it must recover the Newtonian force law in its stated limit and
the free-path/source/clock relations it claims in the relativistic regime.
\end{proposition}
\begin{proof}
H-IF supplies a typed residual derivative; (IF2) supplies its testable physical
momentum-exchange meaning; and (IF3), together with IG1--IG8, supplies a source,
geometry, matter coupling, and observables. F1--F6 establish the flat recovery,
while the remaining listed conditions establish the gravitational recovery and
its empirical domain. Each arrow is a stated premise, so the conclusion applies
only to a realization that proves every map and limit.
\end{proof}

\paragraph{Failure conditions.}
The contract fails as a derivation if it (i) lacks a declared $\mathcal Q_I$,
residual, derivative domain, or response/readout rule; (ii) calls
$-D\mathcal R_I$ a physical newton without proving $K_{\rm N}$; (iii) inserts
$\mathcal F_I$ directly into a stress--energy or metric equation without $L$,
$\Phi$, and compatible exchange accounting; (iv) defines the residual using the
very gravitational field equation, metric, or force law it claims to derive; or
(v) fails the established free/support, tidal, universality, clock, weak-field,
or post-Newtonian targets. A candidate observationally identical to GR plus
standard low-energy matter remains a grounding interpretation unless it states
a finite additional discriminator.

\section{Conditional metric-source recovery of general relativity}\label{sec:gr}
These recovery gates bound the physics companion's
Remark~\ref{PH-prop:sr-gr} and Remark~\ref{PH-cor:reconcile}:
calling both masses change-load establishes neither universal coupling nor
an equivalence-principle prediction.

This section gives the exact classical GR consequence of a supplied metric-source
realization. It does not claim that the generic residual-gradient derivative
selects such a realization. In units $c=1$, let a candidate supply an oriented,
time-orientable four-dimensional Lorentzian manifold $(M,g)$, its Levi--Civita
connection, a non-null boundary prescription, matter fields $\psi$, and a
universal metric coupling $S_{\rm m}[g,\psi]$. For fixed induced boundary metric,
the Dirichlet action is
\[
\begin{aligned}
 S_{\rm GR}&=\frac{1}{2\kappa}\int_M(R-2\Lambda)\sqrt{|g|}\,d^4x
 +\frac{1}{\kappa}\int_{\partial M}\varepsilon K\sqrt{|h|}\,d^3x\\
 &\quad+S_{\rm joint}+S_{\rm m}[g,\psi],
 \qquad
 T_{\mu\nu}:=-\frac{2}{\sqrt{|g|}}\frac{\delta S_{\rm m}}{\delta g^{\mu\nu}},
 \quad \kappa=8\pi G.
\end{aligned}\tag{GR1}
\]
The boundary term and any required joint terms are part of the premise: the bare
Einstein--Hilbert integral does not give a Dirichlet variational problem merely
by fixing $h$. $T_{\mu\nu}$ is a symmetric covariant rank-two source obtained
from matter metric variation, not a scalar information residual.

For allowed variations, stationarity yields
\[
 \boxed{G_{\mu\nu}+\Lambda g_{\mu\nu}=\kappa T_{\mu\nu}}
 \qquad\left(=\frac{8\pi G}{c^4}T_{\mu\nu}\ \text{when $c$ is restored}\right).
 \tag{GR2}
\]
The contracted Bianchi identity is an off-shell geometric identity. Therefore
GR2 requires $\nabla^\mu T_{\mu\nu}=0$ for the total modeled matter source.
Diffeomorphism invariance of $S_{\rm m}$ gives the related Noether identity, and
on the matter equations supplies this conservation statement; components may
exchange energy--momentum and open systems retain boundary flux.

The free-path consequence is separately conditional. With the test-particle
premise $S_{\rm pp}=-m\int ds$, fixed-endpoint worldline variation gives
\[
 \boxed{u^\nu\nabla_\nu u^\mu=0.}\tag{GR3}
\]
A universal matter coupling is required before interpreting this as a
composition-independent free-fall limit. GR2 alone does not make every charged,
spinning, extended, self-gravitating, or quantum body exactly geodesic. In a
weak, quasi-static, asymptotically Minkowskian regime with
$T_{00}\simeq\rho c^2$ and negligible $\Lambda$ at the scale of interest, GR2
and GR3 recover $\nabla^2\Phi=4\pi G\rho$ and
$\ddot{\mathbf x}=-\nabla\Phi$. Clock, light, tidal, and post-Newtonian results
require the corresponding metric solution, universal coupling, and stated
test domain.

\paragraph{The GE selection breakpoint.}
The missing map is not GR2's variation but the noncircular selection of GR1:
\[
 \mathsf{Select}_{\rm GR}:
 (\mathcal F_I,\Hist^W,\Obs,\mathsf{BC},\mathsf{Units})
 \rightrightarrows
 [g,\nabla,S_{\rm GR},S_{\rm m},G,\Lambda]_{\mathrm{Diff}(M)}.
 \tag{GR4}
\]
No such map is currently supplied. It must declare its residual codomain,
admissible geometry/action class, gauge quotient, boundary/initial conditions,
calibration, uncertainty rule, and a uniqueness or empirical selection rule.
Writing the Einstein--Hilbert action, a Lorentzian metric, $G$, $\Lambda$, or
universal coupling directly into its output constraints assumes the result.
Thus GR1--GR3 are an exact conditional GR derivation; GR4 is the remaining
information-gravity derivation problem.

\subsection{Strongest four-dimensional metric-selection result}
The most that an action-selection argument can presently establish is a
conditional classification theorem. Suppose a candidate GR4 selector proves,
rather than inserts, all of the following restrictions: spacetime is fixed to
four dimensions with a declared Lorentzian signature; $g$ is the sole
geometrical gravitational field; the bulk metric equation $E_{\mu\nu}[g]$ is a
local natural function of the metric two-jet, symmetric, and identically
Levi--Civita-divergence-free; no nondynamical background structure occurs; and
the Einstein-tensor coefficient is nonzero. Then the four-dimensional Lovelock
classification gives
\[
 \boxed{E_{\mu\nu}=aG_{\mu\nu}+b g_{\mu\nu}.}\tag{GR5}
\]
With a separately supplied universal metric matter coupling, division by $a$
puts the field equation into GR2 form with $\Lambda=b/a$ and the relative
matter normalization identified as $\kappa=1/a$ in the chosen conventions.
At the action level, the same restricted bulk dynamics is represented by the
Einstein--Hilbert term plus a volume term; boundary completion and
four-dimensional topological/total-derivative additions remain separately
specified.

\paragraph{What GR5 does not select.}
GR5 does not fix $a$, $b$, $G$, $\Lambda$, the metric signature, topology,
boundary conditions, matter species, or universal coupling. It also does not
justify locality, the metric-only premise, absence of extra fields, or the
restriction to second-order metric equations. Relaxing any of those assumptions
admits higher-curvature effective-field-theory terms, nonlocal actions,
metric-affine/torsional theories, extra fields, or other alternatives. A bare
residual-minimization statement cannot supply the restrictions: its domain,
norm, measure, admissible fields, and boundary data already carry the
selection. In particular, a squared equation residual generally yields
higher-order normal equations, while positive reweightings preserve the same
zero set. Therefore GR5 narrows the GE selector to a precise proof obligation;
it is not a derivation of its premises.

\paragraph{Prior-layer realization obstruction.}
No generic residual, entropy, or relation principle currently closes the prior
realization map from GE to the input of GR5. A physical candidate must first
supply: (R1) an event carrier and causal/order relation; (R2) a controlled
manifold-like continuum limit with a declared dimension and time orientation;
(R3) a volume/scale and observable map that determines a Lorentzian metric
rather than a merely conformal or Riemannian structure; and (R4) a dynamical
selection rule suppressing nonlocal, nonmetric, higher-derivative, or
extra-field alternatives. These are not consequences of an untyped residual.
They are the required domain of $\Phi_\epsilon$ and GR4.

Known conditional routes do not remove this obligation. Horizon-thermodynamic
and entanglement-equilibrium arguments begin with Lorentzian causal geometry,
area entropy, and a stress tensor; causal-order reconstruction begins with a
manifold-like causal order and a volume element; information geometry produces
a statistical Riemannian metric; and massless-spin-two consistency begins with
a Lorentzian quantum field theory and a graviton. Each can constrain a
specified gravitational regime, but none selects Lorentzian four-dimensional,
local metric GR and universal matter coupling from a generic information
residual alone. This is an underdetermination result, not an impossibility claim
about a future specified GE microtheory.

\subsection{R1: realized-support proto-causal order}
GE can nevertheless construct a narrower combinatorial object. For a declared
realization trace, let $E^W$ contain typed realized-transition events
\[
 e_t=\langle W,G_t,a_t,G_{t+1},\mathcal R_t,\Obs_t\rangle,
\]
where $a_t$ is an actually selected admissible continuation. Define a primitive
support link $e\leadsto e'$ only when the declared output/support of $e$ is a
necessary input/support of $e'$ at the stated grain. This is a support relation,
not yet a physical influence relation. If the trace supplies (R1a) acyclicity of
$\leadsto$ and (R1b) a finite set of intermediate link-reachable events for
each proposed comparable pair at its declared resolution, define
\[
 e\prec_D e'\quad\Longleftrightarrow\quad
 e\neq e'\ \text{and there is a finite directed $\leadsto$ path from $e$ to $e'$}.
 \tag{R1}
\]
Then $(E^W,\prec_D)$ is a locally finite strict partial order: irreflexivity
follows from R1a, transitivity from path concatenation, and local finiteness
from R1b. It is invariant under relabeling of events and construction order
when all observables depend only on the typed trace and order relation; otherwise
construction-order-sensitive observables must remain declared physical or
operational data.

\paragraph{Promotion gates.}
R1 establishes only a candidate \emph{proto-causal} order. To call
$\prec_D$ physical causal precedence, a realization must additionally prove
that its events are physically individuated rather than states, records, or
propositions; that each link denotes propagating possible influence under a
declared intervention/dynamical model rather than inference, correlation, or
resource support; and that omitted mediators/common causes are controlled. A
Lorentzian-order claim further needs a faithful manifoldlike continuum
realization with dimension and causal regularity. Under the appropriate
continuum hypotheses, causal order fixes only a conformal metric class; a
volume/counting calibration, clocks, and observables are separately needed for
the full metric. The causal-order uniqueness programme for discrete realizations
remains conditional; it cannot be used here as an established geometry theorem.

\section{Conditional electromagnetic alternate-grounding}\label{sec:em}
This is a separate conditional programme, not a claim that electromagnetism is
already present in the inherited vocabulary. It proposes electromagnetic (EM)
physics as a distinct field-edge grounding or physical realization of the same
residual-gradient structure in H-IF: it retains the single core residual
derivative $\mathcal F_I=-D\mathcal R_I$ in (IF1), rather than adding a second
residual force. The proposed EM realization is neither gravity nor a
re-identification of gravitational load-shaped propagation. It has its own
state, gauge, matter, source, exchange, boundary, and observable obligations.

The relevant 01a material does not close those obligations. Its $S^1=U(1)$
discussion motivates a complex amplitude and global phase invariance in a
conditional Born-rule recovery; it supplies no local $U(1)$ connection, field
strength, charge/current, matter coupling, Maxwell equation, or QED
construction. The contingency of a gauge group there is likewise not an EM
derivation.

\begin{hypothesis}[H-EM: electromagnetic field-edge realization]\label{hyp:em}
One candidate supplies, in addition to H-IF, a separately typed electromagnetic
state and clock/readout,
\[
 X_{\rm EM}:\Cont\longrightarrow\State_{\rm EM},\qquad
 \Theta_{\rm EM}:\State_{\rm EM}\longrightarrow
 (\tau_{\rm EM},\mathcal O_{\rm EM}),
\]
with preparation, composition, physical units, and relativistic transformation
rules. It supplies a local $U(1)$ connection $A_\mu$, its curvature/field
strength $F_{\mu\nu}$, and a stated gauge equivalence (for example,
$A_\mu\mapsto A_\mu+\partial_\mu\chi$ with the corresponding transformation of
charged matter). It separately defines a measurable charge/current
$(\rho,j^\mu)$, a matter coupling and physical response law, such as the
Lorentz-force target $m u^\nu\nabla_\nu u^\mu=qF^\mu{}_{\nu}u^\nu$, or an
explicitly different, testable alternative.

The candidate must give an action and equations, or a viable alternative with
equivalent stated conservation and variational/accounting content; gauge-
invariant observables (at minimum a declared operational construction from
$F_{\mu\nu}$, fluxes, and/or holonomies); interface and boundary conditions;
and local and boundary exchange accounting. In a standard covariant realization
this includes a current balance such as $\nabla_\mu j^\mu=0$ and a declared
matter--field stress--energy exchange law. It must state relativistic covariance
and its domain, and must recover Maxwell theory and QED in their respective
tested domains, or name a quantitatively tested alternative with its distinct
predictions and error bounds. None of these requirements is supplied by
H-IF, H-PC, or 01a.
\end{hypothesis}

\subsection{Conditional variational recovery of classical electromagnetism}
A classical electromagnetic recovery can be derived exactly once H-EM supplies
additional geometric and dynamical premises. Let $M$ be an oriented
four-dimensional Lorentzian manifold with metric $g$, let $P\to M$ be a
principal $U(1)$ bundle, and let $A$ be a connection with local curvature
$F=dA$. Let $\Phi$ be charged matter and stipulate the local gauge-invariant
field-edge action
\[
 S_{\rm edge}[A,\Phi]=-
 \frac{1}{4g_{\rm em}^{2}}\int_M F_{\mu\nu}F^{\mu\nu}\sqrt{|g|}\,d^4x
 +S_{\rm m}[g,A,\Phi],\qquad
 J^\nu:=-\frac{1}{\sqrt{|g|}}\frac{\delta S_{\rm m}}{\delta A_\nu}.
 \tag{EM1}
\]
These are new realization premises. In particular, neither the choice of
$U(1)$, the connection, the quadratic Maxwell term, $g_{\rm em}$, nor charged
matter follows from global phase alone or from the generic residual derivative.

For admissible variations of $A$, the action gives
\[
\begin{aligned}
 \delta S_{\rm edge}
 &=\int_M\sqrt{|g|}\left[
 g_{\rm em}^{-2}\nabla_\mu F^{\mu\nu}-J^\nu\right]\delta A_\nu\,d^4x\\
 &\quad-g_{\rm em}^{-2}\int_{\partial M}\sqrt{|h|}\,
 n_\mu F^{\mu\nu}\delta A_\nu\,d^3x.
\end{aligned}\tag{EM2}
\]
With a declared Dirichlet boundary condition (or the corresponding added
Legendre term and declared flux data), stationarity for arbitrary interior
variations therefore yields
\[
 \boxed{\nabla_\mu\bigl(g_{\rm em}^{-2}F^{\mu\nu}\bigr)=J^\nu,
 \qquad dF=0.}\tag{EM3}
\]
The second equation follows from $F=dA$; the first follows from the stipulated
action and source definition. Taking the divergence of the first requires
$\nabla_\nu J^\nu=0$ in the stated smooth classical domain. Equivalently,
local gauge invariance plus the matter equations supplies the corresponding
current identity, with boundary flux retained rather than silently discarded.

Matter response is another separate premise. Given, for example,
\[
 S_{\rm pp}=-m\int ds+q\int A_\mu dx^\mu,\tag{EM4}
\]
worldline variation with fixed endpoints yields
\[
 \boxed{m u^\nu\nabla_\nu u^\mu=qF^\mu{}_{\nu}u^\nu.}\tag{EM5}
\]
Together with (EM3), the electromagnetic stress tensor obeys the matching
field--matter exchange relation only when $S_{\rm m}$ is a
diffeomorphism-covariant matter action, the matter equations hold, and the
specified boundary flux is retained. Thus EM1--EM4 plus those conditions entail
the classical Maxwell field equations, charge balance, and the stated
Lorentz-force response in their declared domain.

\paragraph{Exact breakpoint.}
This is a conditional classical-electrodynamics derivation, not a derivation of
EM1 from GE. A globally defined phase supplies at most a pure-gauge derivative,
not nonzero curvature, a Maxwell kinetic term, a charge representation, or a
physical source. The map from $\mathcal F_I=-D\mathcal R_I$ to the field-edge
action must still be supplied without inserting EM1 into the definition of the
residual. Quantization, photon states, gauge fixing/BRST or an equivalent
physical-state construction, anomaly control, and renormalization are further
QED premises; they do not follow from EM3--EM5.

\paragraph{Non-circularity and falsification.}
H-EM is admissible only if the definition of $\mathcal R_I$, the state/clock
maps, and the translation from the single $\mathcal F_I$ do not assume the
Maxwell or QED action, Maxwell equations, a physical $A_\mu$ or $F_{\mu\nu}$,
charge conservation, the Lorentz-force law, or the experimental EM observables
that the realization claims to recover. It fails as a proposed physical
realization if it lacks separately typed state/clock, local gauge connection or
an explicit viable replacement, charge/current and matter response, an
action/equations and gauge-invariant readout, interface/boundary conditions,
conservation/exchange accounting, relativistic covariance, or the stated
Maxwell/QED-or-alternative recovery. It also fails if a gauge variation changes
an alleged physical observable without a declared symmetry breaking, if the
charge/current or field/matter balance is violated in its claimed domain, or if
its finite alternative prediction is experimentally excluded. These are
falsifiers of H-EM, not of the internal passive-reuse proposition.

This addition does not alter gravity's conditional recovery status. It preserves
the limits on local accelerometry and history-bearing attribution, and 01i
remains unreleased. It supplies neither physical electromagnetism nor Maxwell
theory nor QED; it only makes explicit what a separately implemented EM
grounding would have to establish.

\section{Flat and curved recovery maps}\label{sec:maps}
A candidate realization must present maps rather than metaphors:
\[
\begin{aligned}
 X&:\Cont\longrightarrow\State, &
 P_4&:\State\longrightarrow T^*M,\\
 \Phi&:\Cont\longrightarrow(M,g_{\mu\nu},\nabla,\Obs_{\rm phys}), &
 L&:\Cont\longrightarrow T^{\rm eff}_{\mu\nu}.
\end{aligned}
\]
The domains, codomains, units, transformation rules, approximation errors, and
operational readouts of all four maps are obligations. The project currently
supplies none of them. In particular, a graph is not a Lorentzian manifold, an
alteration cost is not stress--energy, and cycle order is not proper time.

\subsection{Flat-map target}
The following named conditions are required to label the flat equation a
recovery rather than an assumption.
\begin{condition}[F1--F6: flat passive-transport conditions]\label{cond:flat}
A realization must provide: (F1) a flat spacetime structure
$(M,\eta_{\mu\nu},\nabla^{(\eta)})$ with declared signature, physical state
space, and preparation/readout protocol; (F2) a dimensioned covector momentum
map $P_4:\State\to T^*M$, mass-shell convention, and the raised vector
$p^\mu=\eta^{\mu\nu}p_\nu$ with its Lorentz transformation law; (F3) a
proper-time or declared affine-parameter clock map and its timelike/massless
domain; (F4) a free matter action or equations of motion; (F5) the translation
symmetry or explicitly stated balance structure that supports the claimed
conservation; and (F6) an isolation and boundary-flux domain. Only if F1--F6
yield
\[
 \boxed{\frac{dp^\mu}{d\tau}=0}
\]
for the declared free flat regime is this a flat passive-continuation recovery
target met. The equation is not derived in this paper.
\end{condition}

The condition does not confer global momentum conservation on an arbitrary
curved spacetime or on a subsystem. It requires a current, symmetry, or stated
replacement and explicit field/boundary exchange.

\subsection{Load-shaped-map target}
\begin{condition}[IG1--IG8: load-shaped free-propagation conditions]\label{cond:curved}
A realization must provide: (IG1) a controlled map to a Lorentzian metric
$g_{\mu\nu}$ and a declared connection $\nabla$, including whether it is
Levi--Civita or the constraint that determines otherwise; (IG2) a causal and
proper-time clock observable; (IG3) a matter coupling and test-body
approximation; (IG4) a derivation of the metric geodesic law from that coupling;
(IG5) a dimensioned operational load/source map to $T^{\rm eff}_{\mu\nu}$;
(IG6) a metric field equation or explicitly parameterized metric modification;
(IG7) compatible source/exchange and boundary accounting; and (IG8) tidal and
accelerometer observables. Only if IG1--IG8 yield
\[
 \boxed{u^\nu\nabla_\nu u^\mu=0}
\]
in the stated free-body regime is this a load-shaped geodesic recovery target
met. The equation is not derived in this paper.
\end{condition}

A geodesic target alone supplies neither the metric's source nor its dynamics.
For an Einstein-tensor form, the contracted Bianchi identity requires a
compatible source/exchange relation. A candidate must give the tensor type,
units, operational definition, and any $\nabla^\mu T^{\rm eff}_{\mu\nu}=J_\nu$
accounting; it cannot insert an untyped cost into a field equation
\citep{einstein1916,brading2005,szabados2009}.

\section{Equivalence and source obligations}\label{sec:obligations}
\begin{hypothesis}[H-PC: physical passive-continuation bridge]\label{hyp:bridge}
One concrete realization assigns the same declared physical passive-reuse
sector to both the F1--F6 flat regime and the IG1--IG8 load-shaped regime, with
all maps in Section~\ref{sec:maps} defined and measured.
\end{hypothesis}

H-PC is a new hypothesis, not an inherited result. It owes a universal matter
coupling calculation before it can invoke the equivalence principle. That
calculation must state inertial and gravitational response, composition,
binding-energy, spin/history dependence, clock coupling, and the predicted
Eotvos parameter or corresponding observable. Existing tests tightly constrain
such departures \citep{will2014,wagner2012,touboul2022}.

The source obligation is separate. The model must identify a dimensioned source
$T^{\rm eff}_{\mu\nu}$, its exchange law, and metric gravity dynamics. It must
recover or explicitly parameterize deviations from the weak-field targets
\[
 \nabla^2\Phi_N=4\pi G\rho,\qquad \ddot{\mathbf x}=-\nabla\Phi_N
\]
and the applicable redshift, light, clock, and post-Newtonian limits. Neither
universal coupling nor geodesic motion derives Einstein equations;
neither a source equation nor a metaphorical ``load'' derives a geodesic.

\section{Four-way derivation ledger}\label{sec:ledger}
\small\setlength{\tabcolsep}{2pt}\sloppy
\begin{center}
\begin{tabular}{@{}>{\raggedright\arraybackslash}p{0.20\textwidth}>{\raggedright\arraybackslash}p{0.20\textwidth}>{\raggedright\arraybackslash}p{0.24\textwidth}>{\raggedright\arraybackslash}p{0.18\textwidth}@{}}
\toprule
\textbf{Inherited/internal results} & \textbf{New hypotheses} & \textbf{Open physical bridge obligations} & \textbf{Empirical constraints}\\
\midrule
Passive reuse has zero additional alteration cost in its declared class; it is not zero total activity or time. & H-PC: one physical realization has common passive-reuse sectors in both regimes. & Give state identity, physical alteration class, units, support, and operational readout. & Candidate must state which activity, clock, and support measurements distinguish class cost from total expenditure.\\
\addlinespace
Momentum is not force; persistence alone does not prove momentum conservation. & H-F: carried structure maps to measured four-momentum in a flat regime. & Close F1--F6 and derive $dp^\mu/d\tau=0$, including symmetry, isolation, and flux accounting. & Recover inertial transport and measured momentum balance within the declared special-relativistic/\allowbreak Newtonian domain.\\
\addlinespace
Free GR test-body motion has zero proper acceleration; support and tides remain real. & H-G: a dimensioned load maps to controlled geometry and matter coupling. & Close IG1--IG8 and derive $u^\nu\nabla_\nu u^\mu=0$ plus source dynamics. & Meet accelerometer, tidal, weak-field, redshift, light-propagation, and PPN limits.\\
\addlinespace
A present passive continuation is historically underdetermined; proper acceleration is only a local relation. & H-I: a history-bearing attribution record has a specified inference relation. & Define retained changes, source/support history, uncertainty, scale, and comparison protocol without a privileged observer. & Match local accelerometer results while testing momentum/load attribution through the declared historical record.\\
\addlinespace
No project source, metric, connection, equivalence principle, or conservation theorem has been established. & H-U: the coupling is universal in the stated domain. & Calculate composition, binding-energy, spin, and clock effects; give Bianchi-compatible exchange. & Satisfy UFF/EEP and local-Lorentz constraints or state a finite excluded parameter range.\\
\bottomrule
\end{tabular}
\end{center}
\fussy\normalsize\setlength{\tabcolsep}{6pt}

\section{Prior art}\label{sec:priorart}
\begin{sloppypar}
Newton provides the historical inertial/change-law setting; Noether gives the
symmetry--conservation theorem under its action hypotheses \citep{newton1687,noether1918}.
Einstein supplies the metric field theory, while Geroch and Jang formalise a
geodesic principle for suitably small bodies under stated hypotheses
\citep{einstein1916,gerochjang1975}. These sources give targets in their own
formalisms; none maps this project's alteration cost or carried history to
physical momentum or geometry.
\end{sloppypar}

Machian and relational-inertia proposals motivate a comparison but do not close
this bridge \citep{sciama1953,barbourbertotti1977}. Teleparallel gravity gives a
reformulation of already specified GR variables; it is not a license to call
gravity ordinary flat-space inertia, nor a derivation from passive continuation
\citep{aldrovandipereira2013,lisotirioubarrow2011}. Effective field theory
likewise constrains low-energy departures after fields and couplings have been
defined; it does not supply them \citep{donoghue1994,burgess2004}.

\section{Recovery conditions and falsifiers}\label{sec:falsifiers}
The conditional reconciliation is met only if one realization closes F1--F6
and IG1--IG8, supplies H-PC's common implemented dynamics, calculates universal
coupling and source exchange, and recovers the declared tested limits. Even
then it establishes two controlled readings of a law, not a force identity.

A candidate fails as a physical reconciliation if any of the following holds:
\begin{enumerate}[leftmargin=*]
\item its alteration class lacks declared comparison, units or ordering,
support, and a measurement protocol;
\item its free flat regime fails to yield constant measured momentum/four-momentum
under its stated isolation conditions, or it asserts conservation without a
symmetry/current and boundary accounting;
\item its load has no dimensioned source type, or its geometry/source equation
has incompatible divergence or unaccounted exchange;
\item its clock, proper-acceleration, free/support, or tidal observable conflicts
with its claimed regime's established tests;
\item its matter coupling violates UFF, local-Lorentz, redshift, or
post-Newtonian constraints without a surviving quantitative fit;
\item it uses a metric, geodesic principle, universal coupling, conservation
law, or stress tensor as an input while calling it a passive-continuation
derivation; or
\item it attributes a present passive continuation to momentum or gravity
without the declared history-bearing record and inference map; or
\item it has no finite discriminator beyond GR plus standard low-energy matter
while claiming a new physical theory. In that case it is, at most, an
interpretation subject to its own coherence audit.
\end{enumerate}

\section{Conclusion}\label{sec:conclusion}
This paper establishes one internal result and one strict programme condition.
The result is Definition~\ref{def:pc}: reuse of presently operative structure
has zero additional alteration cost in a declared grounding class. The condition
is H-PC: flat inertial transport and load-shaped geodesic motion count as two
physical readings of a common passive-continuation law only after a concrete
realization derives their equations and limits through the named maps and
conditions.

Nothing in the inherited vocabulary derives physical momentum, its conservation,
a metric, connection, source, equivalence principle, geodesic equation, or
quantum-gravity result. Momentum remains distinct from force. GR free fall
remains distinct from a locally felt gravitational force, while support and
tidal effects remain observable. A present in-frame record can establish only
its declared local relation; differentiating momentum from gravity as an
explanatory attribution requires a history-bearing record and the stated
inference map. That requirement supplies no contrary identity claim.

H-EM adds a distinct, conditional electromagnetic field-edge programme while
preserving those limits. It retains one residual-gradient residual derivative
and requires separately typed state/clock, local gauge, charge/current, matter,
dynamics, observables, interface, exchange, covariance, and empirical-recovery
bridges. It does not derive physical electromagnetism, Maxwell theory, or QED;
01i remains unreleased, and gravity recovery remains conditional.

\begin{thebibliography}{99}
\raggedright
\bibitem[Aldrovandi and Pereira(2013)]{aldrovandipereira2013}
R.~Aldrovandi and J.~G.~Pereira, \emph{Teleparallel Gravity: An Introduction}
(Springer, 2013), arXiv:1303.3897.

\bibitem[Barbour and Bertotti(1977)]{barbourbertotti1977}
J.~B.~Barbour and B.~Bertotti, ``Gravity and inertia in a Machian framework,''
\emph{Physical Review D} \textbf{16}, 2738 (1977),
https://doi.org/10.1103/PhysRevD.16.2738.

\bibitem[Bertotti, Iess and Tortora(2003)]{bertotti2003}
B.~Bertotti, L.~Iess and P.~Tortora, ``A test of general relativity using radio
links with the Cassini spacecraft,'' \emph{Nature} \textbf{425}, 374 (2003),
https://doi.org/10.1038/nature01997.

\bibitem[Brading(2005)]{brading2005}
K.~Brading, ``A note on general relativity, energy conservation, and
Noether's theorems,'' arXiv:physics/0509097 (2005).

\bibitem[Burgess(2004)]{burgess2004}
C.~P.~Burgess, ``Quantum gravity in everyday life: General relativity as an
effective field theory,'' \emph{Living Reviews in Relativity} \textbf{7}, 5
(2004), arXiv:hep-th/0311082.

\bibitem[Donoghue(1994)]{donoghue1994}
J.~F.~Donoghue, ``General relativity as an effective field theory,''
\emph{Physical Review D} \textbf{50}, 3874 (1994),
https://doi.org/10.1103/PhysRevD.50.3874.

\bibitem[Einstein(1916)]{einstein1916}
A.~Einstein, ``Die Grundlage der allgemeinen Relativit\"atstheorie,''
\emph{Annalen der Physik} \textbf{49}, 769 (1916),
https://doi.org/10.1002/andp.19163540702.

\bibitem[Geroch and Jang(1975)]{gerochjang1975}
R.~Geroch and P.~S.~Jang, ``Motion of a body in general relativity,''
\emph{Journal of Mathematical Physics} \textbf{16}, 65 (1975),
https://doi.org/10.1063/1.522416.

\bibitem[Li, Sotiriou and Barrow(2011)]{lisotirioubarrow2011}
B.~Li, T.~P.~Sotiriou and J.~D.~Barrow, ``Large-scale structure in $f(T)$
gravity,'' \emph{Physical Review D} \textbf{83}, 104017 (2011),
arXiv:1010.1041.

\bibitem[Newton(1687)]{newton1687}
I.~Newton, \emph{Philosophiae Naturalis Principia Mathematica} (1687),
https://cudl.lib.cam.ac.uk/view/PR-ADV-B-00039-00001/1.

\bibitem[Noether(1918)]{noether1918}
E.~Noether, ``Invariante Variationsprobleme,'' \emph{Nachrichten von der
K\"oniglichen Gesellschaft der Wissenschaften zu G\"ottingen}, 235 (1918);
English translation arXiv:physics/0503066.

\bibitem[Sciama(1953)]{sciama1953}
D.~W.~Sciama, ``On the origin of inertia,'' \emph{Monthly Notices of the Royal
Astronomical Society} \textbf{113}, 34 (1953),
https://doi.org/10.1093/mnras/113.1.34.

\bibitem[Szabados(2009)]{szabados2009}
L.~B.~Szabados, ``Quasi-local energy-momentum and angular momentum in general
relativity,'' \emph{Living Reviews in Relativity} \textbf{12}, 4 (2009),
https://doi.org/10.12942/lrr-2009-4.

\bibitem[Touboul et al.(2022)]{touboul2022}
P.~Touboul \emph{et al.}, ``MICROSCOPE mission: Final results of the test of the
equivalence principle,'' \emph{Physical Review Letters} \textbf{129}, 121102
(2022), https://doi.org/10.1103/PhysRevLett.129.121102.

\bibitem[Wagner et al.(2012)]{wagner2012}
T.~A.~Wagner \emph{et al.}, ``Torsion-balance tests of the weak equivalence
principle,'' \emph{Classical and Quantum Gravity} \textbf{29}, 184002 (2012),
https://doi.org/10.1088/0264-9381/29/18/184002.

\bibitem[Will(2014)]{will2014}
C.~M.~Will, ``The confrontation between general relativity and experiment,''
\emph{Living Reviews in Relativity} \textbf{17}, 4 (2014),
https://doi.org/10.12942/lrr-2014-4.
\end{thebibliography}

\end{document}
